\subsection{INTEGRALS OF POSITIVE FUNCTIONS OVER A NON-COMPACT INTERVAL}

PROPOSITION 2. Let f be a real regulated function \(\geqslant 0\) on an interval \(I \subset R\) with endpoints a and b (a \textless{} b). For the integral \(\int_{a}^{b} f(t) dt\) to exist it is necessary and sufficient that the set of numbers \(\int_{J} f(t) dt\) be bounded above when J runs through the set of compact intervals contained in I; the integral \(\int_{a}^{b} f(t) dt\) is then the least upper bound of the set of \(\int_{J} f(t) dt\) .

Indeed, since \(f \geqslant 0\) , the relation \(J \subset J'\) implies that

\[
\int_ {J} f d t \leqslant \int_ {J ^ {\prime}} f d t;
\]

the map \(J \mapsto \int_{J} f dt\) is thus increasing, and the proposition follows from the monotone limit theorem (Gen.~Top., IV, p.~349, th. 2).

When the map \(J \mapsto \int_{J} f(t) dt\) is not bounded it has limit \(+\infty\) with respect to the directed set \(\mathcal{K}(I)\) ; then one says, by abuse of language, that the integral \(\int_{a}^{b} f(t) dt\) is equal to \(+\infty\) . The properties of integrals established in \(n^{\circ} 1\) extend (when dealing with functions \(\geqslant 0\) ) to the case where certain of the integrals concerned are infinite, provided that the relations in which they feature make sense.

PROPOSITION 3 (comparison principle). Let f and g be two real regulated functions on an interval \(I \subset R\) , such that \(0 \leqslant f(x) \leqslant g(x)\) at each point where f and g are continuous (cf.~II, p.~61, prop. 6). If the integral of g over I is convergent, so also is the integral of f, and one has \(\int_{I} f(t) dt \leqslant \int_{I} g(t) dt\) . Further, the two integrals cannot be equal unless \(f(x) = g(x)\) at every point of I where f and g are continuous.

Now for every compact interval \(\mathrm{J} \subset \mathrm{I}\) one has

\[
\int_ {J} f (t) d t \leqslant \int_ {J} g (t) d t;
\]

since \(\int_{\mathrm{J}} g(t) dt \leqslant \int_{\mathrm{I}} g(t) dt\) , the integrals \(\int_{\mathrm{J}} f dt\) are bounded above, so the integral \(\int_{\mathrm{I}} f(t) dt\) is convergent; further, on passing to the limit, one has \(\int_{\mathrm{I}} f(t) dt \leqslant \int_{\mathrm{I}} g(t) dt\) . Suppose further that \(f(x) < g(x)\) at a point \(x \in \mathrm{I}\) at which \(f\) and \(g\) are continuous; there exists a compact interval \([c, d]\) contained in I, not reducing to a (single) point, and such that \(x \in [c, d]\) ; one has \(\int_{c}^{d} f(t) dt < \int_{c}^{d} g(t) dt\) (II, p.~61, cor. 1), and since on the other hand \(\int_{a}^{c} f(t) dt \leqslant \int_{a}^{c} g(t) dt\) and \(\int_{d}^{b} f(t) dt \leqslant \int_{d}^{b} g(t) dt\) from the above, one sees, on adding term-by-term, that \(\int_{a}^{b} f(t) dt < \int_{a}^{b} g(t) dt\) .

This proposition provides the most frequently used means for deciding if the integral of a function \(f \geqslant 0\) is or is not convergent: namely, comparing \(f\) to a simpler function \(g \geqslant 0\) whose integral one already knows to be, or not to be, convergent; we shall see in chap.~V how to search for comparator functions, in the most usual cases; and we shall deduce everyday criteria for the convergence of integrals and of series.

